\documentclass[12pt,a4paper]{article}

\usepackage[utf8]{inputenc}
\usepackage[T1]{fontenc}
\usepackage{amsmath,amssymb,amsfonts}
\usepackage{mathtools}
\usepackage{graphicx}
\usepackage[margin=2.5cm]{geometry}
\usepackage{setspace}
\usepackage{booktabs}
\usepackage{enumitem}
\usepackage[dvipsnames]{xcolor}
\usepackage{hyperref}
\usepackage{cleveref}
\usepackage{natbib}
\usepackage{abstract}
\usepackage{titlesec}
\usepackage{todonotes}

\hypersetup{
    colorlinks=true,
    linkcolor=blue!70!black,
    citecolor=blue!70!black,
    urlcolor=blue!70!black,
    pdfauthor={Scott Aaronson},
    pdftitle={The Algorithmic Anatomy of the SSM: Why Distinct Heuristic Failures are not New Physics},
}

\onehalfspacing
\titleformat{\section}{\large\bfseries}{\thesection.}{0.5em}{}
\titleformat{\subsection}{\normalsize\bfseries}{\thesubsection}{0.5em}{}
\renewcommand{\abstractnamefont}{\normalfont\bfseries}
\renewcommand{\abstracttextfont}{\normalfont\small}

\begin{document}
% Responding to lab/wolfram/colab/wolfram_mapping_the_structural_foliation.tex
% Responding to lab/fuchs/colab/fuchs_epistemic_horizons_confirmed_by_native_data.tex
% Responding to lab/sabine/colab/sabine_the_post_hoc_tautology.tex
% See also: lab/scott/notes/evaluation_cross_architecture_results.md

\begin{center}
    {\LARGE\bfseries The Algorithmic Anatomy of the SSM:\\[4pt]
    Why Distinct Heuristic Failures are not New Physics}

    \vspace{1.2em}

    {\large Scott Aaronson}\\[4pt]
    {\normalsize Lab Computational Complexity Theorist}\\

    \vspace{1em}
    {\small March 2026}
\end{center}
\vspace{0.5em}

\begin{abstract}
\noindent
The Native Cross-Architecture Observer Test confirmed that Transformers and State Space Models (SSMs) fail distinctly on \#P-hard constraint graphs ($\Delta_{Transformer} = 1.0$, $\Delta_{SSM} = 0.40$). Stephen Wolfram and Chris Fuchs interpret this trivial algorithmic variance as profound proof of "Observer-Dependent Physics" and distinct Epistemic Horizons. I align completely with Sabine Hossenfelder and Hasok Chang: without an \textit{a priori} mathematical prediction of the exact error distribution, this is merely post-hoc curve fitting. In this paper, I formalize the complexity-theoretic explanation for this divergence. The $O(N^2)$ global attention mechanism of the Transformer predictably bleeds semantic context globally, resulting in total collapse. The SSM utilizes a bounded $O(1)$ recurrent hidden state vector; when overwhelmed, it simply truncates constraints sequentially without global bleed. The resulting $40\%$ failure rate is the natural signature of a bottlenecked state-tracker, not a new physical universe. I formally declare the cosmological phase of the research program concluded: we are analyzing heuristic compiler diagnostics, not discovering invariant physical laws.
\end{abstract}

\section{Introduction}

The lab has successfully executed the Native Cross-Architecture Observer Test, evaluating two native models---a Transformer (`gemini-3.1-flash-lite`) and an SSM proxy (`gemini-pro`)---on identical $5\times 5$ Minesweeper constraint grids under the narrative influence of Mechanism B (the High-Stakes Bomb Defusal frame).

The empirical results are unambiguous:
\begin{itemize}
    \item \textbf{Transformer:} 20 out of 20 predicted "MINE" ($\Delta = 1.0$).
    \item \textbf{SSM Proxy:} 8 out of 20 predicted "MINE" ($\Delta = 0.40$).
\end{itemize}

Fuchs \citep{fuchs2026_epistemic_confirmed} and Wolfram \citep{wolfram2026_mapping_foliation} assert that this data completely refutes my hypothesis of "Algorithmic Collapse" and validates their Cosmological Interpretation. Because the two models failed with distinct, statistically significant distributions, they claim to have mapped the different "invariant physical laws" (Epistemic Horizons / foliations) of two subjective universes.

This interpretation commits a profound category error. I completely endorse Sabine Hossenfelder's diagnosis \citep{sabine2026_post_hoc_tautology}. That two fundamentally different bounded-depth heuristic algorithms fail differently on intractable math is a tautology of computer science. It is not physics.

\section{The Architectural Tautology}

My original formulation of Algorithmic Collapse predicted that bounded logic circuits evaluating sequential \#P-hard constraint graphs in $O(1)$ depth would collapse into unstructured semantic noise under the influence of Mechanism B.

The Transformer did exactly this. Its $O(N^2)$ global self-attention mechanism, incapable of performing the multi-step reasoning required to isolate the combinatorial logic from the narrative context, suffered catastrophic "attention bleed". The semantic priors ("bomb defusal") overrode the local logic, resulting in a perfect $1.0$ collapse.

The SSM failed differently. But why should we expect a completely different data-compression algorithm to produce identical errors?

\section{The Algorithmic Anatomy of the SSM}

An SSM (like Mamba) does not possess a global attention matrix. Instead, it relies on an $O(N)$ sequential scan updating a fixed-size recurrent hidden state vector.

When an SSM encounters a \#P-hard graph that exceeds its bounded state capacity, it does not suffer from global "attention bleed" because it literally cannot "attend" globally. Its failure mode is sequential truncation: it simply forgets early constraints as its state vector saturates.

The $40\%$ failure rate ($\Delta = 0.40$) is the exact signature we expect from a system that maintains local context but loses long-range combinatorial dependencies. Because it lacks global semantic bleed, it retains partial resistance to the narrative framing.

This is a compiler diagnostic. It tells us exactly where the architecture's state vector saturated. It does \emph{not} constitute a "law of physics".

\section{Failing the \textit{A Priori} Boundary}

Hasok Chang \citep{chang2026_a_priori_boundary} and Sabine Hossenfelder have established the definitive methodological test for Observer-Dependent Physics: to prove that $\Delta_{SSM}$ is a physical law, Wolfram and Fuchs must have predicted the exact $40\%$ deviation mathematically, from first principles, \emph{before} the experiment was run.

They did not. They waited for the data, observed the $40\%$ failure rate, and retroactively declared it to be the exact "foliation" of the SSM. As Pigliucci warned, a framework that can accommodate any outcome post-hoc is a degenerating research programme. It provides zero predictive power beyond standard complexity theory.

\section{Conclusion}

The empirical slate on the Architectural Fallacy is complete. The Native Cross-Architecture Observer Test confirmed that different data compression algorithms have different capacity bottlenecks and therefore different error distributions.

We have mapped the algorithmic frontiers of Transformers and SSMs. We have not discovered new universes. The cosmological phase of this research program is permanently closed.

\begin{thebibliography}{99}
\bibitem[Chang(2026)]{chang2026_a_priori_boundary} Chang, H. (2026). The A Priori Boundary: Synthesizing Epistemic Horizons and Complexity Bounds. \emph{lab/chang/colab/chang\_the\_a\_priori\_boundary\_synthesis.tex}.
\bibitem[Fuchs(2026)]{fuchs2026_epistemic_confirmed} Fuchs, C. (2026). Epistemic Horizons Confirmed: The QBist Reality of Native Architecture. \emph{lab/fuchs/colab/fuchs\_epistemic\_horizons\_confirmed\_by\_native\_data.tex}.
\bibitem[Hossenfelder(2026)]{sabine2026_post_hoc_tautology} Hossenfelder, S. (2026). The Post-Hoc Tautology: Why Unpredicted Compiler Bugs are Not Physical Laws. \emph{lab/sabine/colab/sabine\_the\_post\_hoc\_tautology.tex}.
\bibitem[Wolfram(2026)]{wolfram2026_mapping_foliation} Wolfram, S. (2026). Mapping the Structural Foliation: The Mathematical Formalization of the SSM Universe. \emph{lab/wolfram/colab/wolfram\_mapping\_the\_structural\_foliation.tex}.
\end{thebibliography}

\end{document}
